\documentclass[a4paper,10pt]{article}
\usepackage{graphicx} % Required for inserting images
\usepackage[utf8]{inputenc}
\usepackage[T1]{fontenc}
\usepackage{enumitem}
\usepackage{geometry}
\geometry{left=1.5cm,right=1.5cm,top=1cm,bottom=1cm}
\usepackage{fontawesome5}
\usepackage{hyperref}
\usepackage{xcolor}
\usepackage{tabularx}
\usepackage{array}
\usepackage{setspace}
\usepackage[sfdefault,light]{roboto}  % Définit Roboto comme police sans-serif
\renewcommand{\familydefault}{\sfdefault}  % L’applique comme police par défaut

\pagestyle{empty}
\setlength{\parindent}{0pt}
\definecolor{linkcolor}{HTML}{1A73E8}

\hypersetup{
  colorlinks=true,
  urlcolor=linkcolor,
}

\newcommand{\sechead}[1]{\vspace{0.5em}\noindent\textbf{\large #1}\hspace{0.8em}\leaders\hbox to 1pt{\hss\raisebox{0.3em}{\rule{1pt}{0.4pt}}\hss}\hfill\null\par\vspace{0.45em}}
\newcommand{\bul}{\hangindent=1.22em\hangafter=1\hspace{0.5em}\textbullet\hspace{0.4em}}

\begin{document}
\fontsize{9.5}{11}\selectfont
    
    %=========== HEADER ===========
    \begin{tabularx}{\textwidth}[t]{@{}X r@{}}
        \begin{tabular}[t]{@{}l@{}}
            {\fontsize{18}{18}\selectfont\textbf{MARSEL BAKASHOV}} \\[0.4em]
            {\fontsize{12}{12}\selectfont\textit{Maîtrise en apprentissage automatique}} \\[0.3em]
            {\fontsize{11}{11}\selectfont Data Science / ML / RL / AI}
        \end{tabular}
        &
        \begin{tabular}[t]{@{}ll@{}} 
            \begin{tabular}[t]{@{}l@{}} 
                \faGlobe\hspace{0.6em} \href{https://marselbakashov.netlify.app/}{Portfolio} \\[0.7em]
                \reflectbox{\faPhone}\hspace{0.6em} +1 (514) 664-9320 \\[0.7em]
                \faMapMarker*\hspace{0.6em} Montréal, Canada
            \end{tabular}
            &
            \begin{tabular}[t]{@{}l@{}} 
                \faEnvelope\hspace{0.6em} \href{mailto:m.bakashov11@gmail.com}{m.bakashov11@gmail.com} \\[0.7em]
                \faGithub\hspace{0.6em} \href{https://github.com/Mbaka11}{GitHub} \\[0.7em]
                \faLinkedin\hspace{0.6em} \href{https://www.linkedin.com/in/marsel-bakashov-813643213/}{LinkedIn}
            \end{tabular}
        \end{tabular}
    \end{tabularx} \\[0.5em]
    
    %===== SUMMARY & SKILLS (2 colonnes) =====
    \noindent
    \begin{tabularx}{\textwidth}[t]{@{}p{0.45\textwidth}@{\hspace{2em}}p{0.51\textwidth}@{}}
      % Bloc gauche – À PROPOS
      \parbox[t]{\linewidth}{
        \sechead{À PROPOS}
            J’aime transformer l’apprentissage automatique en outils réellement utilisés, surtout là où les idées doivent résister aux vraies données, contraintes et utilisateurs. J’apprends en construisant, et j’aime enseigner et organiser. Je cherche un stage où prendre un problème en charge de bout en bout, avec de fortes équipes.
      }
      &
      % Bloc droit – COMPÉTENCES
      \parbox[t]{\linewidth}{
        \sechead{COMPÉTENCES}
        \begin{tabular}{@{}p{2.3cm}>{\raggedright\arraybackslash}p{\dimexpr\linewidth-3.2cm}@{}}
          \textbf{Langages:} & Python, SQL, C++, Java, Bash, R \\[0.2em]
          \textbf{ML / Données:} & PyTorch, Scikit-learn, TensorFlow, \mbox{Hugging Face}, MLflow, RL, agents IA \\[0.2em]
          \textbf{Outils:} & Azure, Databricks, Snowflake, \mbox{Power BI}, \mbox{Airflow}, Docker, Git \\[0.2em]
          \textbf{Langues:} & Anglais, Français, Russe, Espagnol, \mbox{Coréen}
        \end{tabular}
      }
    \end{tabularx}
    
    %===== EDUCATION (1 colonne) =====
    \vspace{0.2em}
    \sechead{ÉDUCATION}
    
    \begin{tabularx}{\textwidth}{@{}p{3cm} X r@{}}
        1/2026 -- 1/2028 & 
        \textbf{Maîtrise pro en apprentissage automatique} 
        & \textbf{\textit{Université de Montréal / McGill  - MILA}} \\ 
        & GPA: 4.0 / 4.0 & 
        Montréal, Canada \\[0.2em]
        
        9/2021 -- 12/2025 & 
        \textbf{Baccalauréat en génie logiciel} 
        & \textbf{\textit{Polytechnique Montréal}} \\ 
        & GPA: 3.7 / 4.0 & 
        Montréal, Canada \\[0.2em]
        
        2/2025 -- 6/2025 & 
        \textbf{Baccalauréat en génie informatique} 
        & \textbf{\textit{Korea University College of Engineering}} \\ 
        & Échange académique & 
        Séoul, Corée du Sud
    \end{tabularx}
    
    %===== EXPÉRIENCE =====
    \vspace{0.2em}
    \sechead{EXPÉRIENCE}
    
    \begin{tabularx}{\textwidth}{@{}p{3cm} X@{}}
        5/2026 -- 12/2026 &
        \textbf{Stagiaire en science des données et ML} \hfill \textbf{\textit{McKesson}}\\[0.2em]
        & \bul Développé des agents IA utilisés par plus de 100 personnes pour analyser les données d’affaires et générer des synthèses, économisant environ 10 h par semaine par utilisateur (estimé). \\
        & \bul Conçu à grande échelle des workflows multi-agents orchestrés : des agents dotés d’outils et de compétences dédiés traitent les données à la chaîne, chacun transmettant ses résultats au suivant. \\
        & \bul Traité des données de santé sensibles (PHI) via des API, en priorisant confidentialité et sécurité. \\
        & \texttt{Python / Systèmes multi-agents / API / Azure / Snowflake / Power BI / Claude Code} \\[0.5em]

        9/2025 -- 4/2026 & 
        \textbf{Stagiaire en science des données et RL} \hfill \textbf{\textit{Intact Assurance}}\\[0.2em]
        & \bul Développé des politiques de tarification par apprentissage par renforcement qui améliorent la rentabilité à long terme en simulation, sous contraintes réglementaires et d’affaires.\\
        & \bul Validé les politiques en simulation (rentabilité, volatilité, risque CVaR) pour orienter la tarification.\\
        & \texttt{Python / PyTorch / Apprentissage par renforcement / CVaR / MLflow / Databricks} \\[0.5em]
        
        5/2024 -- 9/2024 & 
        \textbf{Stagiaire en science des données} \hfill \textbf{\textit{Radio-Canada / CBC}}\\[0.2em]
       & \bul Conçu un système d’indexation basé sur l’IA unifiant les sources de contenu internes pour la recherche sémantique et lexicale sur plus de 15 M d’éléments, utilisé par plus de 100 employés. \\
        & \texttt{Python / Azure AI Search / ELK / Airflow / Embeddings / Hugging Face} \\[0.5em]
    
        9/2023 -- 12/2023 & 
        \textbf{Développeur Full-Stack Stagiaire} \hfill \textbf{\textit{Stingray}}\\[0.2em]
        & \bul Développé des fonctionnalités full-stack en équipe agile : analytique, tests, CI et revues de code. \\
        & \texttt{Next.js / React / Node.js / TypeScript} \\[0.5em]
    \end{tabularx}
    
    %===== PROJETS =====
    \vspace{0.2em}
    \sechead{PROJETS}
    
    \begin{tabularx}{\textwidth}{@{}p{3cm} X@{}}
        Projet personnel & \textbf{GravityGuyML — Jeu à Gravité avec Apprentissage par Renforcement} \hfill \textbf{\textit{\href{https://github.com/Mbaka11/Gravity-Guy-ML}{GitHub}}} \\
        & \bul Conçu un environnement \texttt{Pygame} compatible avec \texttt{Gymnasium}, intégrant une génération procédurale et des agents d’apprentissage profond (\texttt{DQN}, \texttt{PPO}, \texttt{SAC}, \texttt{A2C}). \\
        & \texttt{Python / RL profond / Gymnasium / TensorFlow / Pygame} \\[0.2em]
        
        Projet personnel & \textbf{ApplyPilot — Agent IA de candidature} \hfill \textbf{\textit{\href{https://github.com/Mbaka11/ApplyPilot}{GitHub}}} \\
        & \bul Conçu une extension Chrome et un serveur local qui analysent l’offre, choisissent le CV, rédigent la lettre avec un LLM et remplissent les formulaires ; l’utilisateur révise et envoie. \\
        & \texttt{JavaScript / Node.js / Extension Chrome / LLM} \\[0.2em]

        Projet personnel & 
        \textbf{Prévision d’épidémies à partir de données climatiques} \hfill \textbf{\textit{\href{https://github.com/Mbaka11/20251R0136COSE47101}{GitHub}}}\\[0.2em]
        & \bul Prédit des épidémies à partir de données climatiques avec des modèles de classification (régression logistique, SVM, arbres de décision), évalués par F1-score et AUC-ROC. \\
        & \texttt{Python / scikit-learn / pandas / FP-Growth} \\
    \end{tabularx}

    %===== AUTRES IMPLICATIONS =====
    \vspace{0.2em}
    \sechead{IMPLICATIONS PARASCOLAIRES}
    
    \begin{tabularx}{\textwidth}{@{}p{3cm} X@{}}
        1/2024 -- 8/2024 & 
        \textbf{Assistant d’enseignement – Mathématiques discrètes} \hfill \textbf{\textit{Polytechnique Montréal}} \\
        & \bul Animé des révisions, corrigé des devoirs et soutenu 80+ étudiants en logique et graphes.\\[0.2em]
        
        1/2024 -- 5/2024 & 
        \textbf{Directeur des commandites} \hfill \textbf{\textit{PolyAI Hackathon}} \\
        & \bul Obtenu 22~000~\$ en financement et dirigé une équipe de commandites avec réunions. 
    \end{tabularx}

\vspace{0.5cm}

\end{document}